
\documentclass[12pt,a4paper]{article}

% Поддержка русского языка
\usepackage[T2A]{fontenc}
\usepackage[utf8]{inputenc}
\usepackage[russian]{babel}

% Математика
\usepackage{amsmath,amssymb,amsthm}

% Настройка страницы
\usepackage[a4paper,margin=2.5cm]{geometry}

\begin{document}
	
	\tableofcontents
	
	\newpage
	\section{Обозначения}
	Фиксируем $x \in \mathbb{R}$\\
	Пусть $f(x) = (1, x, \ldots, x^m)^\mathrm{T}$, $c=f'(x)$\\
	
	\noindent Рассмотрим план
	\[
	\xi = 
	\left(
	\begin{matrix}
		x_1 \ x_2 \ldots x_n \\
		\omega_1 \ \omega_2 \ldots \omega_n	
	\end{matrix}
	\right)
	\]
	
	\noindent Та самая система для плана с $n$ узлами имеет вид:
	\[
	c = h M(\xi) p = h \sum_{k=1}^{n} \omega_k <p,f(x_k)> f(x_k)
	\]	
	
	\noindent или, c учётом $h = (\lambda ^*)^{-1}$ и чередования знаков т.е. $<p,f(x_k)> = (-1)^{k+1}$
	\[
	\lambda ^*
	\begin{pmatrix}
		0 \\
		1 \\
		\vdots \\
		m x^{m-1}
	\end{pmatrix} = 
	\sum_{k=1}^{n} a_k 
	\begin{pmatrix}
		1 \\
		x_k \\
		\vdots \\
		x_k^{m}
	\end{pmatrix}, \quad a_k = \omega_k <p,f(x_k)>
	\]
	
	\noindent Введём линейные операторы $A_1, A_2, A: \mathbb{P}_m \to \mathbb{R}$:
	\[
	A_1(p) := \lambda^* p'(x)
	\]
	\[
	A_2(p) := \sum_{k=1}^{n} a_k p(x_k)
	\]
	\[
	A(p) := A_1(p) - A_2(p)
	\]
	
	\noindent Нетрудно видеть, что наша система эквивалентна вот такому:
	\[
	\lbrace 1, x, \ldots, x^m \rbrace \subset \ker A
	\]
	
	\noindent Т.к $A$ линейно, то это в свою очередь эквивалентно
	\begin{equation}
		\mathbb{P}_m \subset \ker A
		\label{eq:cond1}
	\end{equation}
	
	
	\newpage
	\section{Чебышёвские планы}
	\subsection{Ищем веса в чебышёвских планах}
	
	
	\noindent (\textbf{в этом разделе} $\mathbf{n = m+1}$)\\ 
	
	
	\noindent Вспоминаем, что \eqref{eq:lagr_expansion}  $\forall p \in \mathbb{P}_m$ даёт нам
	\[
	p(t) = 
	\sum_{k=1}^{m+1} p(x_k) L_k(t), \  \forall t \in \mathbb{R}
	\]
	
	\noindent Дифференцируя это равенство по $t$ и подставляя $t=x$ получаем 
	\[
	p'(x) = 
	\sum_{k=1}^{m+1} p(x_k) L_k'(x)
	\]
	
	
	\noindent Получили, что $\forall p \in \mathbb{P}_m$
	\begin{equation}
		\lambda^* p'(x) = 
		\sum_{k=1}^{m+1} p(x_k) ( \lambda^* L_k'(x) ) 
		\label{eq:re1_2}
	\end{equation}
	
	\noindent Но из условия \eqref{eq:cond1} 
	\[
	A(p) = 0, \quad \forall p \in \mathbb{P}_m
	\]
	
	\noindent То есть $\forall p \in \mathbb{P}_m$
	\begin{equation}
		\lambda^* p'(x) = 
		\sum_{k=1}^{m+1} p(x_k) ( a_k ) 
		\label{eq:re2_2}
	\end{equation}
	
	\noindent Из \eqref{eq:re1_2} и \eqref{eq:re2_2} получаем $\forall p \in \mathbb{P}_m$
	\[
	\sum_{k=1}^{m+1} p(x_k) ( a_k ) = \sum_{k=1}^{m+1} p(x_k) ( \lambda^* L_k'(x) ) 
	\] 
	
	\noindent Выбирая в качестве $p$ штуки которые равны $1$ и $0$ где надо ($L_k(x) ???$), делаем вывод
	\[
	a_k = \lambda^* L_k'(x), \ \forall k \in 1 \ldots m+1 
	\]
	\[
	\omega_k (-1)^{k+1} = \lambda^* L_k'(x), \ \forall k \in 1 \ldots m+1 
	\]
	\[
	\omega_k  = (-1)^{k+1} \lambda^* L_k'(x), \ \forall k \in 1 \ldots m+1 
	\]
	
	
	\noindent Юзая условие на веса $\sum_{k=1}^{m+1} \omega_k = 1$ ищем лямбду:
	\[
	1 = \sum_{k=1}^{m+1} \omega_k = \sum_{k=1}^{m+1} (-1)^{k+1} \lambda^* L_k'(x) = \lambda^* \sum_{k=1}^{m+1} (-1)^{k+1}  L_k'(x)
	\]
	
	\noindent получаем
	\[
	\lambda^* = \left( \sum_{k=1}^{m+1} (-1)^{k+1}  L_k'(x) \right) ^ {-1}
	\]
	
	\noindent финальное выражение для весов
	\[
	\omega_k  = 
	\frac
	{
		(-1)^{k+1} L_k'(x)
	}
	{
		\sum_{i=1}^{m+1} (-1)^{i+1}  L_i'(x)
	}, \ \forall k \in 1 \ldots m+1 
	\]
	
	\[
	L_k(x) = \underset{k \neq j}{\prod_{j=1}^{m+1}} \frac{x - x_j}{x_k - x_j}, \ k \in 1 \ldots m+1
	\]
	
	\subsection{Ограничение на $x$ от весов}
	
	Мы требуем $\forall k \in 1 \ldots m+1 \ \omega_k > 0 $
	
	\noindent То есть нужно решить систему неравенств относительно $x$:
	\[
	\omega_k  = (-1)^{k+1} \lambda^* L_k'(x)  > 0 , \ \forall k \in 1 \ldots m+1 
	\]
	
	
	
	\subsection{Ищем узлы в чебышёвских планах}
	
	\noindent А узлы это просто точки альтернанса полинома Чебышёва 1-го рода 
	\[
	T_{m}(x) = \cos( m \arccos(x))
	\]
	\noindent Они имеют вид 
	\[
	x_{k+1} = \cos \left( \frac{(m-k) \pi}{m} \right), \ k \in 0 \ldots m
	\]
	
	\noindent Для $m=4$ многочлен Чебышёва выглядит вот так:
	\[
	T_{4}(x) = \cos( 4 \arccos(x)) = 8x^4 - 8x^2 + 1
	\]
	
	\noindent А его точки альтернанса:
	\begin{align*}
		x_1 &= \cos \left( \frac{4\pi}{4} \right) = -1\\
		x_2 &= \cos \left( \frac{3\pi}{4} \right) = -\frac{\sqrt{2}}{2}\\
		x_3 &= \cos \left( \frac{2\pi}{4} \right) = 0\\
		x_4 &= \cos \left( \frac{1\pi}{4} \right) = \frac{\sqrt{2}}{2}\\
		x_5 &= \cos \left( \frac{0\pi}{4} \right) = 1\\
	\end{align*}
		
		
	\noindent Теперь найдем веса, для этого сначала выпишем базисные многочлены и их производные
	
	\begin{align*}
		L_1(x)
		&=
		\frac{
			\left(x+\frac{\sqrt{2}}{2}\right)
			x
			\left(x-\frac{\sqrt{2}}{2}\right)
			(x-1)
		}{
			\left(-1+\frac{\sqrt{2}}{2}\right)
			(-1)
			\left(-1-\frac{\sqrt{2}}{2}\right)
			(-2)
		}
		\\
		&=
		\left(x+\frac{\sqrt{2}}{2}\right)
		x
		\left(x-\frac{\sqrt{2}}{2}\right)
		(x-1)
		\\
		&=
		x^4-x^3-\frac{1}{2}x^2+\frac{1}{2}x
		\\[4mm]
		L'_1(x) &= 4x^3 -3x^2-x+\frac{1}{2}
	\end{align*}
	
	
	\begin{align*}
		L_2(x)
		&=
		\frac{
			(x+1)
			x
			\left(x-\frac{\sqrt{2}}{2}\right)
			(x-1)
		}{
			\left(-\frac{\sqrt{2}}{2}+1\right)
			\left(-\frac{\sqrt{2}}{2}\right)
			(-\sqrt{2})
			\left(-\frac{\sqrt{2}}{2}-1\right)
		}
		\\
		&=
		-2(x+1)
		x
		\left(x-\frac{\sqrt{2}}{2}\right)
		(x-1)
		\\
		&=
		-2x^4+\sqrt{2}x^3+2x^2-\sqrt{2}x \\[4mm]
		L'_2(x) &= -8x^3 + 3\sqrt{2}x^2 + 4x - \sqrt{2} 	
	\end{align*}
		
	\begin{align*}
		L_3(x)
		&=
		\frac{
			(x+1)
			\left(x+\frac{\sqrt{2}}{2}\right)
			\left(x-\frac{\sqrt{2}}{2}\right)
			(x-1)
		}{
			(1)
			\left(\frac{\sqrt{2}}{2}\right)
			\left(-\frac{\sqrt{2}}{2}\right)
			(-1)
		}
		\\
		&=
		2(x+1)
		\left(x+\frac{\sqrt{2}}{2}\right)
		\left(x-\frac{\sqrt{2}}{2}\right)
		(x-1)
		\\
		&=
		2x^4-3x^2+1,
		\\[4mm]
		L'_3(x) &= 8x^3 - 6x
	\end{align*}
	
	\begin{align*}
		L_4(x)
		&=
		\frac{
			(x+1)
			\left(x+\frac{\sqrt{2}}{2}\right)
			x
			(x-1)
		}{
			\left(1+\frac{\sqrt{2}}{2}\right)
			\sqrt{2}
			\left(\frac{\sqrt{2}}{2}\right)
			\left(-1+\frac{\sqrt{2}}{2}\right)
		}
		\\
		&=
		-2(x+1)
		\left(x+\frac{\sqrt{2}}{2}\right)
		x
		(x-1)
		\\
		&=
		-2x^4-\sqrt{2}x^3+2x^2+\sqrt{2}x,
		\\[4mm]
		L'_4(x) &= -8x^3 - 3\sqrt{2}x^2 + 4x + \sqrt{2}
	\end{align*}
		
	\begin{align*}
		L_5(x)
		&=
		\frac{
			(x+1)
			\left(x+\frac{\sqrt{2}}{2}\right)
			x
			\left(x-\frac{\sqrt{2}}{2}\right)
		}{
			(2)
			\left(1+\frac{\sqrt{2}}{2}\right)
			(1)
			\left(1-\frac{\sqrt{2}}{2}\right)
		}
		\\
		&=
		(x+1)
		\left(x+\frac{\sqrt{2}}{2}\right)
		x
		\left(x-\frac{\sqrt{2}}{2}\right)
		\\
		&=
		x^4+x^3-\frac{1}{2}x^2-\frac{1}{2}x\\[4mm]
		L'_5(x) &= 4x^3 + 3x^2 -x - \frac{1}{2}
	\end{align*}
	
	
	
	\begin{align*}
		\sum_{k=1}^{5}(-1)^{k+1}L_k'(x)
		&=
		L_1'(x)-L_2'(x)+L_3'(x)-L_4'(x)+L_5'(x) = \\
		&=
		(4x^3 -3x^2-x+\frac{1}{2})\\
		&
		- (-8x^3 + 3\sqrt{2}x^2 + 4x - \sqrt{2})\\
		&
		+ (8x^3 - 6x)\\
		&
		-(-8x^3 - 3\sqrt{2}x^2 + 4x + \sqrt{2})\\
		&
		+(4x^3 + 3x^2 -x - \frac{1}{2}) =\\
		& =
		32x^3-16x \\
		& =
		16x(2x^2-1)
	\end{align*}
	
	\noindent Отсюда лямбда равна
	\[
	\lambda^* = 
	\left( \sum_{k=1}^{m+1} (-1)^{k+1}  L_k'(x) \right) ^ {-1} = 
	\frac{1}{16x(2x^2-1)}
	\]
	
	\noindent Теперь находим веса из формулы
	\[\omega_k  = (-1)^{k+1} \lambda^* L_k'(x), \ \forall k \in 1 \ldots m+1 \]
	\begin{align*}
		\omega_1 &= 
		(-1)^{1+1} \lambda^* L_1'(x) = 
		\frac{(4x^3 -3x^2-x+\frac{1}{2})}{16x(2x^2-1)}\\[3mm]
		\omega_2 &= 
		(-1)^{2+1} \lambda^* L_2'(x) = 
		\frac{- (-8x^3 + 3\sqrt{2}x^2 + 4x - \sqrt{2})}{16x(2x^2-1)}\\[3mm]
		\omega_3 &= 
		(-1)^{3+1} \lambda^* L_3'(x) = 
		\frac{(8x^3 - 6x)}{16x(2x^2-1)}\\[3mm]
		\omega_4 &= 
		(-1)^{4+1} \lambda^* L_4'(x) = 
		\frac{-(-8x^3 - 3\sqrt{2}x^2 + 4x + \sqrt{2})}{16x(2x^2-1)}\\[3mm]
		\omega_5 &= 
		(-1)^{5+1} \lambda^* L_5'(x) = 
		\frac{(4x^3 + 3x^2 -x - \frac{1}{2})}{16x(2x^2-1)}\\
	\end{align*}
	
	\noindent В конце концов получаем что план имеет вид
	
	\[
	\begin{pmatrix}
	-1 & -\frac{\sqrt2}{2} & 0 & \frac{\sqrt2}{2} & 1 \\
	\frac{(4x^3 -3x^2-x+\frac{1}{2})}{16x(2x^2-1)} &
	\frac{- (-8x^3 + 3\sqrt{2}x^2 + 4x - \sqrt{2})}{16x(2x^2-1)} &
	\frac{(8x^3 - 6x)}{16x(2x^2-1)} &
	\frac{-(-8x^3 - 3\sqrt{2}x^2 + 4x + \sqrt{2})}{16x(2x^2-1)} &
	\frac{(4x^3 + 3x^2 -x - \frac{1}{2})}{16x(2x^2-1)}
	\end{pmatrix}
	\]
	
	\newpage
	\section{Получебышёвские планы}
	\subsection{Ищем веса в получебышёвских планах}
	\noindent (\textbf{А теперь} $\mathbf{n = m}$)\\ 
	Теперь рассмотрим 
	\[
	p_0(x) = \prod_{k=1}^{m} (x-x_k)
	\]
	
	\noindent Для него
	\[A(p_0) = 0\]
	
	\noindent И, также 
	\[A_2(p_0) = 0\]
	
	\noindent Отсюда 
	\[A_1(p_0) = \lambda^* p_0'(x) =  0\]
	
	\noindent При фиксированном $x$ получаем условие на $x_k, k \in 1 \ldots m$ (лямбда не ноль)
	\[
	p_0'(x) = \frac{d}{dt} \left[\prod_{k=1}^{m} (t - x_k) \right]_{t=x} = \ 0
	\]
	
	\noindent Вспоминаем, что \eqref{eq:lagr_expansion}  $\forall p \in \mathbb{P}_m$ даёт нам
	\[
	p(t) = 
	\sum_{k=1}^{m} p(x_k) L_k(t) + 
	a_m(p) \prod_{k=1}^{m} (t-x_k) = 
	\sum_{k=1}^{m} p(x_k) L_k(t) + 
	a_m(p) p_0(t)
	\]
	
	\noindent Дифференцируя это равенство по $t$ получаем 
	\[
	p'(t) = 
	\sum_{k=1}^{m} p(x_k) L_k'(t) + 
	a_m(p) p_0'(t)
	\]
	
	\noindent Подставляя $t=x$ убиваем остаток
	\[
	p'(x) = 
	\sum_{k=1}^{m} p(x_k) L_k'(x)
	\]
	
	\noindent Получили, что $\forall p \in \mathbb{P}_m$
	\begin{equation}
		\lambda^* p'(x) = 
		\sum_{k=1}^{m} p(x_k) ( \lambda^* L_k'(x) ) 
		\label{eq:re1}
	\end{equation}
	
	\noindent Но из условия \eqref{eq:cond1} 
	\[
	A(p) = 0, \quad \forall p \in \mathbb{P}_m
	\]
	
	\noindent То есть $\forall p \in \mathbb{P}_m$
	\begin{equation}
		\lambda^* p'(x) = 
		\sum_{k=1}^{m} p(x_k) ( a_k ) 
		\label{eq:re2}
	\end{equation}
	
	\noindent Из \eqref{eq:re1} и \eqref{eq:re2} получаем $\forall p \in \mathbb{P}_m$
	\[
	\sum_{k=1}^{m} p(x_k) ( a_k ) = \sum_{k=1}^{m} p(x_k) ( \lambda^* L_k'(x) ) 
	\] 
	
	\noindent Выбирая в качестве $p$ штуки которые равны $1$ и $0$ где надо, делаем вывод
	\[
	a_k = \lambda^* L_k'(x), \ \forall k \in 1 \ldots m 
	\]
	\[
	\omega_k (-1)^{k+1} = \lambda^* L_k'(x), \ \forall k \in 1 \ldots m 
	\]
	\[
	\omega_k  = (-1)^{k+1} \lambda^* L_k'(x), \ \forall k \in 1 \ldots m 
	\]
	
	
	\noindent Юзая условие на веса $\sum_{k=1}^{m} \omega_k = 1$ ищем лямбду:
	\[
	1 = \sum_{k=1}^{m} \omega_k = \sum_{k=1}^{m} (-1)^{k+1} \lambda^* L_k'(x) = \lambda^* \sum_{k=1}^{m} (-1)^{k+1}  L_k'(x)
	\]
	
	\noindent получаем
	\[
	\lambda^* = \left( \sum_{k=1}^{m} (-1)^{k+1}  L_k'(x) \right) ^ {-1}
	\]
	
	\noindent финальное выражение для весов
	\[
	\omega_k  = 
	\frac
	{
		(-1)^{k+1} L_k'(x)
	}
	{
		\sum_{i=1}^{m} (-1)^{i+1}  L_i'(x)
	}, \ \forall k \in 1 \ldots m 
	\]
	
	\[
	L_k(x) = \underset{k \neq j}{\prod_{j=1}^{m}} \frac{x - x_j}{x_k - x_j}, \ k \in 1 \ldots m
	\]
	
	\subsection{Ограничение на $x$ от весов}
	
	Мы требуем $\forall k \in 1 \ldots m \ \omega_k > 0 $
	
	\noindent То есть нужно решить систему неравенств относительно $x$:
	\[
	\omega_k  = (-1)^{k+1} \lambda^* L_k'(x) > 0 , \ \forall k \in 1 \ldots m 
	\]
	
	\newpage
	\subsection{А теперь ищем узлы}
	
	\subsection*{Если растяжка вправо}
	
	\noindent В таком сетапе считаем
	\[
	-1 = x_1 < x_2 < \ldots < x_m < 1 < z
	\]
	
	\noindent где $x_k = a t_{k-1} + b, k \in 1, \ldots m$ \\
	и $z=a t_m + b$ (её игнорим)\\
	для неких $a, b \in \mathbb{R}$, где $a > 0$ (мы их не знаем, но они явно и не нужны, главное что какое-то линейное преобразование) \\
	а $t_k, \  k \in 0, \ldots, m$ есть точки альтернанса полинома Чебышёва 1-го рода 
	\[
	T_{m}(x) = \cos( m \arccos(x))
	\]
	\noindent Они имеют вид 
	\[
	t_k = \cos \left( \frac{ (m-k) \pi}{m} \right), \ k \in 0 \ldots m
	\]
	
	\noindent Уравнения описывающие растяжение (отношение длин интервалов между точками альтернанса сохраняется)
	\begin{align*}
		\frac{x_3 - x_2}{x_2 - x_1} &= \frac{t_3 - t_2}{t_2 - t_1}\\ 
		\frac{x_4 - x_3}{x_2 - x_1} &= \frac{t_4 - t_3}{t_2 - t_1}\\
		&\cdots\\
		\frac{x_m - x_{m-1}}{x_2 - x_1} &= \frac{t_m - t_{m-1}}{t_2 - t_1}\\
	\end{align*}
	\noindent уже известно, что $x_1 = -1$, значит из первого уравнения $x_3$ линейно выражается через $x_2$.\\
	из второго $x_4$ линейно выражается через $x_2$ и $x_3$ то есть линейно через $x_2$\\
	в конце концов все $x_k, \ k \in 1 \ldots m $ линейно выражаются через $x_2$\\
	Для окончательного нахождения плана достаточно найти $x_2$ 
	\textbf{ИЛИ} найти $a$ и $b$.\\
	
	
	\noindent Заметим, что узлы удовлетворяют
	\[
	p_0'(x) = 
	\frac{d}{dt} \left[\prod_{k=1}^{m} (t - x_k) \right]_{t=x} = 
	\frac{d}{dt} \left[\prod_{k=1}^{m} (t - (at_{k-1} + b) ) \right]_{t=x} = 
	0
	\]
	
	\noindent Введём новую переменную $s = \frac{t-b}{a}$
	\[
	\frac{d}{dt} \left[\prod_{k=1}^{m} (t - (at_{k-1} + b) ) \right] =
	\frac{1}{a}\frac{d}{dt} \left[\prod_{k=1}^{m} \left( \frac{t-b}{a} - t_{k-1} \right) \right] 
	\]
	\[
	\frac{d}{dt} \left[\prod_{k=1}^{m} \left( \frac{t-b}{a} - t_{k-1} \right) \right] = 0 
	\quad \Longleftrightarrow \quad
	\frac{d}{ds} \left[\prod_{k=1}^{m} \left( s - t_{k-1} \right) \right] = 0
	\]
	
	\noindent В $x_k = a t_{k-1} + b, k \in 1, \ldots m$ подставляем $k=1$ и получаем
	\[
	x_1 = a t_0 + b
	\]
	\[
	-1 = a (-1) + b
	\]
	\[
	a = b + 1
	\]
	
	\noindent Когда мы найдем значения $s$ из
	\[
	\frac{d}{ds} \left[\prod_{k=1}^{m} \left( s - t_{k-1} \right) \right] = 
	\frac{d}{ds} \left[\prod_{k=0}^{m-1} \left( s - t_{k} \right) \right] = 0
	\]
	
	\noindent Мы воспользуемся тем, что
	\[
	s |_{t=x} = \frac{x-b}{a} = \frac{x-b}{b+1}
	\]
	
	\noindent И отсюда найдём $b$, затем $a$ \\
	после этого можно явно выписать выражения для $x_k \  k \in 1 \ldots m$
	
	\noindent Если обратиться к полиномам Чебышёва 2-го рода то можно переписать уравнение на $s$ в виде:
	\[
	\frac{d}{ds} \left[ (s-t_0) \prod_{k=1}^{m-1} \left( s - t_{k} \right) \right] = 
	\frac{d}{ds} \left[ (s+1) \prod_{k=1}^{m-1} \left( s - t_{k} \right) \right] = 
	C \frac{d}{ds} \left[ (s+1) U_{m-1}(s) \right] = 0
	\]
	
	\noindent Где
	\[
	U_{m-1}(s) = \frac{\sin(m \arccos(s))}{\sin(\arccos(s))}
	\]
	
	\noindent Положим $m=4$ и в лоб решим уравнение на $s$
	\[
	\frac{d}{ds} \left[\prod_{k=0}^{m-1} \left( s - t_{k} \right) \right] = 0
	\]
	
	\noindent Явно выпишем то что дифференцируется
	\begin{align*}
	&\prod_{k=0}^{m-1} \left( s - t_{k} \right) = \\
	&= (s + 1) \left(s + \frac{1}{\sqrt2} \right) (s) \left(s - \frac{1}{\sqrt2} \right) = \\
	&= (s^2 + s) \left( s^2 - \frac{1}{2} \right) \\
	&= s^4 + s^3 - \frac{1}{2}s^2 - \frac{1}{2}s 
	\end{align*}
	
	\begin{align*}
	&\frac{d}{ds} \left[\prod_{k=0}^{m-1} \left( s - t_{k} \right) \right] = \\ 
	&= \frac{d}{ds} \left[ s^4 + s^3 - \frac{1}{2}s^2 - \frac{1}{2}s \right] = \\
	&= 4s^3 + 3s^2 - s - \frac{1}{2} = 0
	\end{align*}
	
	\newpage
	
	\noindent Пытаемся решить кубическое...\\

	Разделим уравнение на $4$:
	
	$$
	s^3+\frac34s^2-\frac14s-\frac18=0.
	$$
	
	Для приведения кубического уравнения к нормальной форме сделаем замену
	
	$$
	s=y-\frac14.
	$$
	
	Подставляя это выражение в исходное уравнение, получаем
	
	$$
	\left(y-\frac14\right)^3
	+\frac34\left(y-\frac14\right)^2
	-\frac14\left(y-\frac14\right)
	-\frac18=0.
	$$
	
	Раскрывая скобки, имеем
	
	$$
	y^3-\frac34y^2+\frac3{16}y-\frac1{64}
	+\frac34\left(y^2-\frac12y+\frac1{16}\right)
	-\frac14y+\frac1{16}
	-\frac18=0.
	$$
	
	После приведения подобных слагаемых квадратичный член сокращается:
	
	$$
	y^3-\frac7{16}y-\frac1{32}=0.
	$$
	
	Таким образом, получили нормальную форму кубического уравнения
	
	$$
	\boxed{
		y^3-\frac7{16}y-\frac1{32}=0.
	}
	$$
	
	Для решения воспользуемся тригонометрической формой решения приведённого кубического уравнения
	
	$$
	y^3+py+q=0.
	$$
	
	В нашем случае
	
	$$
	p=-\frac7{16},
	\qquad
	q=-\frac1{32}.
	$$
	
	Поскольку $p<0$, положим
	
	$$
	y=2\sqrt{-\frac p3}\cos\theta.
	$$
	
	Имеем
	
	$$
	2\sqrt{-\frac p3}
	=
	2\sqrt{\frac7{48}}
	=
	\sqrt{\frac7{12}},
	$$
	
	поэтому
	
	$$
	y=\sqrt{\frac7{12}}\cos\theta.
	$$
	
	После подстановки в уравнение $y^3+py+q=0$ и использования формулы
	
	$$
	4\cos^3\theta-3\cos\theta=\cos(3\theta)
	$$
	
	получаем
	
	$$
	\cos(3\theta)
	=
	\frac{3q}{2p}\sqrt{-\frac3p}.
	$$
	
	Вычислим правую часть:
	
	$$
	\frac{3q}{2p}
	=
	\frac{3(-1/32)}
	{2(-7/16)}
	=
	\frac3{28},
	$$
	
	а
	
	$$
	\sqrt{-\frac3p}
	=
	\sqrt{\frac{48}{7}}
	=
	4\sqrt{\frac37}.
	$$
	
	Следовательно,
	
	$$
	\cos(3\theta)
	=
	\frac3{28}\,4\sqrt{\frac37}
	=
	\frac3{7}\sqrt{\frac37}
	=
	\frac{3\sqrt{21}}{49}.
	$$
	
	Таким образом,
	
	$$
	\boxed{
		\cos(3\theta)=\frac{3\sqrt{21}}{49}.
	}
	$$
	
	Обозначим
	
	$$
	\varphi=
	\arccos\left(\frac{3\sqrt{21}}{49}\right).
	$$
	
	Тогда три значения угла имеют вид
	
	$$
	\theta_k=\frac{\varphi+2\pi k}{3},
	\qquad
	k=0,1,2.
	$$
	
	Следовательно, три корня уравнения для $y$ равны
	
	$$
	y_k=
	\sqrt{\frac7{12}}
	\cos\left(
	\frac13
	\arccos\left(\frac{3\sqrt{21}}{49}\right)
	+\frac{2\pi k}{3}
	\right),
	\qquad
	k=0,1,2.
	$$
	
	Возвращаясь к переменной $s$, поскольку
	
	$$
	s=y-\frac14,
	$$
	
	получаем
	
	$$
	\boxed{
		s_k=
		-\frac14+
		\sqrt{\frac7{12}}
		\cos\left(
		\frac13
		\arccos\left(\frac{3\sqrt{21}}{49}\right)
		+\frac{2\pi k}{3}
		\right),
		\qquad
		k=0,1,2.
	}
	$$
	
	Численные значения корней:
	
	$$
	\boxed{
		s_1 \approx -0.8723,
		\qquad
		s_2 \approx -0.3222,
		\qquad
		s_0 \approx 0.4446.
	}
	$$
	
	
	
	\noindent Выражаем параметры растяжения через $s$
	\begin{align*}
	s |_{t=x} &= \frac{x-b}{a} = \frac{x-b}{b+1}\\
	s (b+1) &= x-b\\
	b &= \frac{x-s}{s+1}\\
	a &= b + 1 = \frac{x-s}{s+1} + 1 = \frac{x+1}{s+1}
	\end{align*}
	
	\noindent Узлы получаем через
	\[
	x_k = a t_{k-1} + b = \frac{(x+1)t_{k-1} + x-s}{s+1} =
	\frac{ x(1 + t_{k-1}) + (t_{k-1} - s) }{s+1} , \qquad \ k \in 1, \ldots m
	\]
	
	\subsection{Ограничение на $x$ от узлов}
	
	\noindent Вспоминаем геометрию узлов:
	\[
	-1 = x_1 < x_2 < \ldots < x_m < 1 < z
	\]
	
	\noindent Отсюда, с учетом растяжения вправо ограничение на $x$:
	\[
	x_m = x_m(x) \in (t_{m-1}, \ 1)
	\]
	
	\newpage
	\subsection*{Если растяжка влево}
	
	\noindent В таком сетапе считаем
	\[
	z < -1 < x_1 < \ldots < x_{m-1} < x_m = 1
	\]
	
	\noindent где $x_k = a t_{k} + b, k \in 1, \ldots m$ \\
	и $z=a t_0 + b$, $a > 0$ (её игнорим)\\
	для неких $a, b \in \mathbb{R}$ (мы их не знаем, но они явно и не нужны, главное что какое-то линейное преобразование) \\
	а $t_k, \  k \in 0, \ldots, m$ есть точки альтернанса полинома Чебышёва 1-го рода 
	\[
	T_{m}(x) = \cos( m \arccos(x))
	\]
	\noindent Они имеют вид 
	\[
	t_k = \cos \left( \frac{ (m-k) \pi}{m} \right), \ k \in 0 \ldots m
	\]
	
	\noindent Уравнения описывающие растяжение (отношение длин интервалов между точками альтернанса сохраняется)
	\begin{align*}
		\frac{x_{m-1} - x_{m-2}}{x_m - x_{m-1}} &= \frac{t_{m-1} - t_{m-2}}{t_m - t_{m-1}}\\[0.1in]
		\frac{x_{m-2} - x_{m-3}}{x_m - x_{m-1}} &= \frac{t_{m-1} - t_{m-2}}{t_m - t_{m-1}}\\[0.1in]
		&\cdots\\[0.1in]
		\frac{x_2 - x_1}{x_m - x_{m-1}} &= \frac{t_2 - t_1}{t_m - t_{m-1}}\\
	\end{align*}
	
	\noindent уже известно, что $x_m = -1$, значит из первого уравнения $x_{m-2}$ линейно выражается через $x_{m-1}$.\\
	из второго $x_{m-3}$ линейно выражается через $x_{m-2}$ и $x_{m-1}$ то есть линейно через $x_{m-1}$\\
	в конце концов все $x_k, \ k \in 1 \ldots m $ линейно выражаются через $x_{m-1}$\\
	Для окончательного нахождения плана достаточно найти $x_{m-1}$ \textbf{ИЛИ} найти $a$ и $b$.\\
	
	
	\noindent Для поиска узлов используем
	\[
	p_0'(x) = 
	\frac{d}{dt} \left[\prod_{k=1}^{m} (t - x_k) \right]_{t=x} = 
	\frac{d}{dt} \left[\prod_{k=1}^{m} (t - (at_{k} + b) ) \right]_{t=x} = 
	0
	\]
	
	\noindent Введём новую переменную $s = \frac{t-b}{a}$
	\[
	\frac{d}{dt} \left[\prod_{k=1}^{m} (t - (at_{k} + b) ) \right] =
	\frac{1}{a}\frac{d}{dt} \left[\prod_{k=1}^{m} \left( \frac{t-b}{a} - t_{k} \right) \right] 
	\]
	\[
	\frac{d}{dt} \left[\prod_{k=1}^{m} \left( \frac{t-b}{a} - t_{k} \right) \right] = 0 
	\quad \Longleftrightarrow \quad
	\frac{d}{ds} \left[\prod_{k=1}^{m} \left( s - t_{k} \right) \right] = 0
	\]
	
	\noindent В $x_k = a t_{k} + b, k \in 1, \ldots m$ подставляем $k=m$ и получаем
	\[
	x_m = a t_m + b
	\]
	\[
	1 = a (1) + b
	\]
	\[
	a = 1 - b
	\]
	
	\noindent Когда мы найдем значения $s$ из
	\[
	\frac{d}{ds} \left[\prod_{k=1}^{m} \left( s - t_{k} \right) \right] = 0
	\]
	
	\noindent Мы воспользуемся тем, что
	\[
	s |_{t=x} = \frac{x-b}{a} = \frac{x-b}{1 - b}
	\]
	
	\noindent И отсюда найдём $b$, затем $a$ \\
	после этого можно явно выписать выражения для $x_k \  k \in 1 \ldots m$
	
	
	\noindent Если обратиться к полиномам Чебышёва 2-го рода то можно переписать уравнение на $s$ в виде:
	\[
	\frac{d}{ds} \left[ (s-t_m) \prod_{k=1}^{m-1} \left( s - t_{k} \right) \right] = 
	\frac{d}{ds} \left[ (s-1) \prod_{k=1}^{m-1} \left( s - t_{k} \right) \right] = 
	C \frac{d}{ds} \left[ (s-1) U_{m-1}(s) \right] = 0
	\]
	
	\noindent Где
	\[
	U_{m-1}(s) = \frac{\sin(m \arccos(s))}{\sin(\arccos(s))}
	\]
	
	\noindent Положим $m=4$ и в лоб решим уравнение на $s$
	\[
	\frac{d}{ds} \left[\prod_{k=1}^{m} \left( s - t_{k} \right) \right] = 0
	\]
	
	\noindent Явно выпишем то что дифференцируется
	\begin{align*}
		&\prod_{k=1}^{m} \left( s - t_{k} \right) = \\
		&=  \left(s + \frac{1}{\sqrt2} \right) (s)  \left(s - \frac{1}{\sqrt2} \right)  (s - 1) = \\
		&= (s^2 - s) \left( s^2 - \frac{1}{2} \right) \\
		&= s^4 - s^3 - \frac{1}{2}s^2 + \frac{1}{2}s
	\end{align*}
	
	\begin{align*}
		&\frac{d}{ds} \left[\prod_{k=1}^{m} \left( s - t_{k} \right) \right] = \\ 
		&= \frac{d}{ds} \left[ s^4 - s^3 - \frac{1}{2}s^2 + \frac{1}{2}s \right] = \\
		&= 4s^3 - 3s^2 - s + \frac{1}{2} = 0
	\end{align*}
	
	\newpage
	
	\noindent Пытаемся решить кубическое...\\
	
	Разделим уравнение на $4$:
	
	$$
	s^3-\frac34s^2-\frac14s+\frac18=0.
	$$
	
	Для приведения кубического уравнения к нормальной форме сделаем замену
	
	$$
	s=y+\frac14.
	$$
	
	Подставляя это выражение в исходное уравнение, получаем
	
	$$
	\left(y+\frac14\right)^3
	-\frac34\left(y+\frac14\right)^2
	-\frac14\left(y+\frac14\right)
	+\frac18=0.
	$$
	
	Раскрывая скобки, имеем
	
	$$
	y^3+\frac34y^2+\frac3{16}y+\frac1{64}
	-\frac34\left(y^2+\frac12y+\frac1{16}\right)
	-\frac14y-\frac1{16}
	+\frac18=0.
	$$
	
	После приведения подобных слагаемых квадратичный член сокращается:
	
	$$
	y^3-\frac7{16}y+\frac1{32}=0.
	$$
	
	Таким образом, получили нормальную форму кубического уравнения
	
	$$
	\boxed{
		y^3-\frac7{16}y+\frac1{32}=0.
	}
	$$
	
	Для решения воспользуемся тригонометрической формой решения приведённого кубического уравнения
	
	$$
	y^3+py+q=0.
	$$
	
	В нашем случае
	
	$$
	p=-\frac7{16},
	\qquad
	q=\frac1{32}.
	$$
	
	Поскольку $p<0$, положим
	
	$$
	y=2\sqrt{-\frac p3}\cos\theta.
	$$
	
	Имеем
	
	$$
	2\sqrt{-\frac p3}
	=
	2\sqrt{\frac7{48}}
	=
	\sqrt{\frac7{12}},
	$$
	
	поэтому
	
	$$
	y=\sqrt{\frac7{12}}\cos\theta.
	$$
	
	После подстановки в уравнение $y^3+py+q=0$ и использования формулы
	
	$$
	4\cos^3\theta-3\cos\theta=\cos(3\theta)
	$$
	
	получаем
	
	$$
	\cos(3\theta)
	=
	\frac{3q}{2p}\sqrt{-\frac3p}.
	$$
	
	Вычислим правую часть:
	
	$$
	\frac{3q}{2p}
	=
	\frac{3(1/32)}
	{2(-7/16)}
	=
	-\frac3{28},
	$$
	
	а
	
	$$
	\sqrt{-\frac3p}
	=
	\sqrt{\frac{48}{7}}
	=
	4\sqrt{\frac37}.
	$$
	
	Следовательно,
	
	$$
	\cos(3\theta)
	=
	-\frac3{28}\,4\sqrt{\frac37}
	=
	-\frac3{7}\sqrt{\frac37}
	=
	-\frac{3\sqrt{21}}{49}.
	$$
	
	Таким образом,
	
	$$
	\boxed{
		\cos(3\theta)=-\frac{3\sqrt{21}}{49}.
	}
	$$
	
	Обозначим
	
	$$
	\varphi=
	\arccos\left(-\frac{3\sqrt{21}}{49}\right).
	$$
	
	Тогда три значения угла имеют вид
	
	$$
	\theta_k=\frac{\varphi+2\pi k}{3},
	\qquad
	k=0,1,2.
	$$
	
	Следовательно, три корня уравнения для $y$ равны
	
	$$
	y_k=
	\sqrt{\frac7{12}}
	\cos\left(
	\frac13
	\arccos\left(-\frac{3\sqrt{21}}{49}\right)
	+\frac{2\pi k}{3}
	\right),
	\qquad
	k=0,1,2.
	$$
	
	Возвращаясь к переменной $s$, поскольку
	
	$$
	s=y+\frac14,
	$$
	
	получаем
	
	$$
	\boxed{
		s_k=
		\frac14+
		\sqrt{\frac7{12}}
		\cos\left(
		\frac13
		\arccos\left(-\frac{3\sqrt{21}}{49}\right)
		+\frac{2\pi k}{3}
		\right),
		\qquad
		k=0,1,2.
	}
	$$
	
	Численные значения корней:
	
	$$
	\boxed{
		s_0 \approx 0.8723,
		\qquad
		s_1 \approx 0.3222,
		\qquad
		s_2 \approx -0.4446.
	}
	$$
	
	\noindent Выражаем параметры растяжения через $s$
	
	\begin{align*}
		s |_{t=x} &= \frac{x-b}{a} = \frac{x-b}{1-b}\\
		s (1-b) &= x-b\\
		b &= \frac{x-s}{1-s}\\
		a &= 1-b = 1-\frac{x-s}{1-s}
		= \frac{1-x}{1-s}
	\end{align*}
	
	\noindent Узлы получаем через
	\[
	x_k = a t_k + b =
	\frac{(1-x)t_k + x-s}{1-s} =
	\frac{ x(1 - t_k)  + (t_k - s)}{1-s}  ,
	\qquad
	k \in 1, \ldots m
	\]
	
	
	
	\subsection{Ограничение на $x$ от узлов}
	
	\noindent Вспоминаем геометрию узлов:
	\[
	z < -1 < x_1 < \ldots < x_{m-1} < x_m = 1
	\]
	
	\noindent Отсюда, с учетом растяжения влево ограничение на $x$:
	\[
	x_1 = x_1(x) \in (-1, \ t_1)
	\]
	
	
	\newpage
	\section{Два узла на границе}
	
	\subsection{Веса}
	
	\noindent Они получаются такими же как и в получебышёвском случае:
	\[
	\omega_k  = 
	\frac
	{
		(-1)^{k+1} L_k'(x)
	}
	{
		\sum_{i=1}^{m} (-1)^{i+1}  L_i'(x)
	}, \ \forall k \in 1 \ldots m 
	\]
	
	\[
	L_k(x) = \underset{k \neq j}{\prod_{j=1}^{m}} \frac{x - x_j}{x_k - x_j}, \ k \in 1 \ldots m
	\]
	
	\subsection{Ограничение на $x$ от весов}
	
	Мы требуем $\forall k \in 1 \ldots m \ \omega_k > 0 $
	
	\noindent То есть нужно решить систему неравенств относительно $x$:
	\[
	\omega_k  = 
	(-1)^{k+1} \lambda^* L_k'(x)> 0 , \ \forall k \in 1 \ldots m 
	\]
	
	\subsection{Узлы}
	
	Теперь узлы должны быть расположены следующим образом:
	\[
	-1 = x_1 < x_2 < \ldots < x_{m-1} < x_m = 1
	\]
	
\end{document}